\section{Two-block uncle Hamiltonians as actual parent limits}
\label{sec:two_block_uncle}

The uncle construction takes a punctured perturbation limit of the canonical
parent Hamiltonian. The binding source is
\cite[continuity-of-projector lemma, span-invariance lemma, and form theorem]{FernandezGonzalez2015Uncle};
these are in \texttt{Papers/1210.6613/UncleHMPS.tex}, lines 314--340,
890--902, and 912--953.
The word ``gauge'' in the span lemma allows an invertible linear map on
virtual matrices. In particular, multiplying both off-diagonal blocks by
one scalar need not be a similarity transformation.

\begin{lemma}[Local continuity at an injective boundary map]
    \label{lem:uncle_local_projector_continuity}
    \lean{ContinuousLinearMap.continuousAt_range_starProjection_of_injective}
    \lean{MPSTensor.continuousAt_groundSpaceES_starProjection_of_isNBlkInjective}
    \lean{MPSTensor.continuousAt_parentInteractionES_of_isNBlkInjective}
    \leanok
    \uses{def:ground_space_map_es,def:parent_interaction}
    Let $T_x$ be a family of finite-dimensional boundary maps continuous at
    $x_0$. If $T_{x_0}$ is injective, its range projector is continuous at
    $x_0$. Consequently, fix a window length $L$ and a tensor family
    continuous at $x_0$. If the products of exactly $L$ tensor matrices at
    $x_0$ span the full virtual matrix algebra, the canonical parent
    interaction on $L$ sites is continuous at that point.
\end{lemma}
\begin{proof}
    \leanok
    Injectivity is open. In a neighborhood of $x_0$, write the projector as
    $T_x(T_x^\dagger T_x)^{-1}T_x^\dagger$ using the existing inverse-Gram
    formula. Inversion is continuous at the invertible base-point Gram
    operator. This is the continuous-basis argument of the source, with
    inverse Gram matrices in place of explicit Gram--Schmidt coordinates.
\end{proof}

\begin{lemma}[Invertible off-diagonal virtual rescaling]
    \label{lem:uncle_virtual_rescaling}
    \lean{MPSTensor.twoBlockVirtualScale}
    \lean{MPSTensor.trace_twoBlockVirtualScale_mul}
    \lean{MPSTensor.twoBlockVirtualScale_inv}
    \lean{MPSTensor.groundSpaceMap_twoBlockVirtualScale_one}
    \lean{MPSTensor.groundSpace_twoBlockVirtualScale_one}
    \leanok
    \uses{def:ground_space_map,def:ground_space_parent}
    For $z\ne0$, the linear map
    $\mathcal L_z\bigl(\begin{smallmatrix}X&Y\\Z&W\end{smallmatrix}\bigr)
        =\bigl(\begin{smallmatrix}X&zY\\zZ&W\end{smallmatrix}\bigr)$
    is invertible and preserves the one-site physical span of a tensor.
\end{lemma}
\begin{proof}
    \leanok
    Its inverse is $\mathcal L_{z^{-1}}$. The bilinear trace identity
    $\tr(\mathcal L_z(M)X)=\tr(M\mathcal L_z(X))$ transfers the rescaling to
    the arbitrary boundary matrix. Surjectivity of this boundary change
    gives equality of the physical ranges. No conjugation or similarity
    assumption is used.
\end{proof}

\begin{definition}[Perturbed blocks and the regular uncle tensor]
    \label{def:two_block_uncle_tensor}
    \lean{MPSTensor.uncleBlockMatrix}
    \lean{MPSTensor.twoBlockPerturbation}
    \lean{MPSTensor.uncleTensorInterpolant}
    \lean{MPSTensor.uncleTensor}
    \lean{MPSTensor.uncleTensorInterpolant_zero}
    \leanok
    Let $A_i,B_i$ be square virtual blocks, and let $P_i,Q_i,R_i,L_i$ be
    arbitrary perturbation blocks of the matching sizes. Define
    \begin{align}
        C_i(\varepsilon) & =
        \begin{pmatrix}A_i+\varepsilon P_i & \varepsilon R_i     \\
               \varepsilon L_i     & B_i+\varepsilon Q_i\end{pmatrix},
        \label{eq:uncle_perturbation} \\
        U_{ij}           & =
        \begin{pmatrix}A_iA_j        & A_iR_j+R_iB_j \\
               L_iA_j+B_iL_j & B_iB_j\end{pmatrix}.
        \label{eq:uncle_explicit_tensor}
    \end{align}
    The regular interpolant $U(\varepsilon)$ has diagonal blocks
    $(A_i+\varepsilon P_i)(A_j+\varepsilon P_j)+\varepsilon^2R_iL_j$
    and $(B_i+\varepsilon Q_i)(B_j+\varepsilon Q_j)+\varepsilon^2L_iR_j$;
    its off-diagonal blocks are those of $U$ plus respectively
    $\varepsilon(P_iR_j+R_iQ_j)$ and
    $\varepsilon(L_iP_j+Q_iL_j)$. In particular $U(0)=U$.
\end{definition}

\begin{lemma}[Regularization of the blocked perturbation]
    \label{lem:uncle_regularization}
    \lean{MPSTensor.twoBlockVirtualScale_blockTensor}
    \lean{MPSTensor.continuous_uncleTensorInterpolant}
    \leanok
    \uses{def:two_block_uncle_tensor,lem:uncle_virtual_rescaling}
    The interpolant is continuous, and for every $\varepsilon\ne0$,
    $U(\varepsilon)=\mathcal L_{\varepsilon^{-1}}
        (C(\varepsilon)\mathbin{\circ}C(\varepsilon))$.
\end{lemma}
\begin{proof}
    \leanok
    Multiply the four matrix blocks and cancel the nonzero factor
    $\varepsilon$ in each off-diagonal block. All remaining entries are
    polynomials in $\varepsilon$.
\end{proof}

\begin{theorem}[Canonical two-block uncle interaction and Hamiltonian]
    \label{thm:two_block_uncle_parent_limit}
    \lean{MPSTensor.uncleInteractionES}
    \lean{MPSTensor.uncleHamiltonianES}
    \lean{MPSTensor.tendsto_parentInteractionES_twoBlockPerturbation}
    \lean{MPSTensor.tendsto_parentHamiltonianES_twoBlockPerturbation}
    \lean{MPSTensor.tendsto_parentHamiltonianES_twoBlockPerturbation_real}
    \leanok
    \uses{lem:uncle_local_projector_continuity,lem:uncle_regularization,
        def:blocked_configuration_isometry,def:parent_hamiltonian}
    Suppose the explicit uncle tensor $U$ is injective. The actual canonical
    two-site parent interactions of $C(\varepsilon)$ converge, in operator
    norm as $\varepsilon\to0$ through nonzero complex parameters, to
    $h'_P=\Pi[U]$ on the original two physical sites. At every fixed periodic
    length $N\ge2$, the actual parent Hamiltonians converge in operator norm
    to $H'_P=\sum_{i=1}^N(h'_P)_{i,i+1}$. The Hamiltonian limit also holds along real
    perturbation parameters.
\end{theorem}
\begin{proof}
    \leanok
    \uses{lem:uncle_local_projector_continuity,lem:uncle_virtual_rescaling,
        lem:uncle_regularization}
    For nonzero parameters, virtual rescaling preserves the physical range
    and hence the canonical projector. The regular interpolant has injective
    value $U$ at zero, so its projector converges to $\Pi[U]$. The existing
    physical-blocking isometry returns this projector to the original
    two-site Hilbert space. Each translated interaction is expressed through
    the existing cyclic restriction and its adjoint. Finite sums and these
    fixed operator compositions preserve operator-norm convergence.
\end{proof}

The source assumes $C=A\oplus B$ is block-injective. The local limit proof
needs only the stated injectivity of $U$, so the formal result also covers
other block data satisfying that condition as an algebraic operator-limit
statement, without additional claims about their standard-form physical
interpretation. Arbitrary diagonal perturbations
$P,Q$ have the same limit for fixed $A,B,R,L$. No assertion about generic
perturbations, an algebraic exceptional set, gaplessness, or an
infinite-volume Hamiltonian is included here.
